\scanpage{84}
% Broad pencil strokes cancel this entire face. The smaller cancellations are local.
\deleted{\uline{Théorème} 3.1. Soient $S$ un \deleted{pré}schéma affine, $\deleted{G}\handwritten{H}$ un $S$-groupe de type multiplicatif, $F$ un faisceau quasi-cohérent sur $S$, sur lequel $G$ opère (Exp 1,
4.7.) Alors \deleted{pour tout $i\quad0$, la cohomologie de Hochschild} on a}
\[
\deleted{H^i(\deleted{G}\handwritten{H},F)=0\qquad\text{pour }i\quad0,}
\]
\deleted{où $H^{\unsure{\bullet|\prime|i}}$ est la cohomologie de Hochschild étudiée dans Exp.1 5.}

\deleted{En effet, d'après loc cit , 5.3. , la cohomologie de Hochschild
s'explicite, en\deleted{\unsure{t}} termes des\deleted{\unsure{it}} anneaux affines $A$ de $S$ et $B$ de $G$, et du
\deleted{\unsure{mo}} module $M$ sur $A$ définissant $F$, comme cohomologie d'un complexe
\deleted{$C^{\unsure{n|p|*}}(G,F)$, dont la} de $A$-modules $C^{\unsure{\circ|\bullet}}(\deleted{G}\handwritten{H},F)$ , \uline{dont la formation commute au
changement de base} $A\quad A'$ . Par suite, pour un changement de base
$A\quad A'$ avec $A'$ \uline{plat} sur $A$, on aura}
\[
\deleted{H^i(\deleted{G}\handwritten{H}',F')=H^i(\deleted{G}\handwritten{\unsure{H|G}},F)\otimes_A A',}
\]
\deleted{et par suite, si on suppose même $A'$ fidèlement plat sur $A$, pour prouver
que $H^i(\deleted{G}\handwritten{H},F)$ est nul, il suffit de prouver qu'il en est ainsi de $H^i(G',F')$.
Cette remarque nous ramène à vérifier 3.1. dans le cas où $G$ est diagonalisable, dans ce cas cela a été prouvé dans Exp 1 5.3.3.}

\deleted{Utilisant les résultats de Exp.III , nous allons tirer de 3.1.
diverses conséquences de nature géométrique:}

\deleted{\uline{Théorème} 3.2. Soient $S$ un \deleted{pré}schéma\deleted{, $G$ un groupe de typ} affine, $S_{\unsure{o|0}}$
un sous-schéma affine défini par un Idéal $J$ \uline{de carré nul} , $\deleted{G}\handwritten{H}$ un groupe
de type multiplicatif sur $S$, $\deleted{H}\handwritten{G}$ un \deleted{\unsure{g}}préschéma en groupes quelconque sur $S$,
$u,v:\deleted{G}\handwritten{H}\quad\deleted{H}\handwritten{G}$ deux homomorphismes de groupes tels que $u_o=v_o$ , où
$u_o,v_o\mathrel{\unsure{\mathbin{:}|=}}\deleted{G}\handwritten{H}_o\quad\deleted{H}\handwritten{G}_o$ sont les homomorphismes déduits de $u,v$ par changement de
base $S_o\quad S$. Alors il existe un $g$ \deleted{\unsure{$G$}} $\in$ \inserted{$G(S)$} tel que $v=\operatorname{int}(g)\ u$ \handwritten{et $g_o=e$.}}

\deleted{Cela résulte de Exp III \qquad et de 3.1.}
% archive sheet number: 83
